\begin{pastexam}[label={px:nchu-113-dsa-16}]{中興資工所甲組 113 學年度 選擇題第 16 題}
（單選；答錯不倒扣）

Given an array of integers, which of the following sorting algorithms would be the most efficient in terms
of time complexity for sorting an almost sorted array where only two elements are swapped? (5\%)
\begin{enumerate}[label=(\Alph*)]
\item Bubble Sort
\item Quick Sort
\item Merge Sort
\item Heap Sort
\end{enumerate}
\end{pastexam}

\begin{pastsoln}
本書採用的答案是 (A)。

出題的意思是：「幾乎排好」的輸入，有提早結束檢查的氣泡排序（某一趟完全沒有交換就停，資結書第~14.2~節）
只要很少幾趟。 例如對調的是相鄰兩個元素，第一趟把它們換回來，第二趟沒有交換就停，總共 $\Theta(n)$，
比其他三個都快：
\begin{itemize}
\item (B) 快速排序取第一個或最後一個當基準時，幾乎排好的輸入正是最壞情況，$\Theta(n^2)$（資結書第~14.5~節）；
隨機基準也還是期望 $\Theta(n \log n)$。
\item (C) 合併排序不管輸入，都是 $\Theta(n \log n)$。
\item (D) 堆積排序先建堆積就把順序打亂，仍是 $\Theta(n \log n)$。
\end{itemize}

要注意，「只有兩個元素對調」若兩者相距很遠，氣泡排序不一定快。 設位置 $i < j$ 的兩個數對調：大的數在
第 $i$ 格，一趟就往右換到第 $j$ 格；小的數在第 $j$ 格，每一趟只往左移一格，要 $j - i$ 趟才到位，總共
$\Theta(n(j - i))$。 例如 $1$ 和 $n$ 對調，要 $\Theta(n^2)$，反而比合併排序慢。 真正對這種輸入都好的是
插入排序：它的時間是 $\Theta(n + \text{逆序數})$，兩個數對調只產生 $2(j - i) - 1$ 個逆序，總共 $\Theta(n)$；
但插入排序不在選項裡。 照題目的意思（對調的兩個數離得不遠），選 (A)。
\end{pastsoln}
